% Part II, sections that do not depend on the results (included by main_en.tex).

\section{One backbone for the policy and the judge}\label{sec:shared}

\kev{}~\cite{kev4b} is a LoRA adapter of rank 16 and a pointer head on Qwen3.5-4B-Base (revision \texttt{1001bb4d}).
The policy we train is a LoRA adapter on the same base. A single copy of the base weights can therefore serve both
models. It carries two adapters: \texttt{default}, the policy (rank 32, $\alpha=64$, on every attention, MLP and Gated
DeltaNet projection), and \texttt{kev}, \kev{}'s released adapter, frozen and applied unmerged. Only the policy adapter
is trained, so the base weights never change and \kev{} reads exactly the backbone it was trained on. Scoring a
completion switches the active adapter to \texttt{kev} and reads \kev{}'s pointer head with its calibration temperature.
The judge then sees the same backbone and adapter weights as a separately loaded \kev{}. Its outputs differ only by
merged versus unmerged LoRA arithmetic in bf16, which we measure (\S\ref{sec:res-shared}).

The alternative is to load \kev{} as a second model. We measure the GPU memory of both designs in the same process
with the same training step: allocated memory after loading, and peak memory over a policy update on a fixed batch
followed by \kev{} scoring.

\section{GRPO setup}\label{sec:grpo}

\paragraph{Policy and data.} The policy starts from Qwen3.5-4B-Base with a zero-initialised LoRA, so step 0 is the base
model. Training problems are MATH~\cite{hendrycks2021math} training problems of levels 3--5. A fixed random 500 of them
(seed 12345) are held out as the validation set, which is the only data used for model selection. The other
training problems are visited in a seeded random order. None of the MATH-500 test problems is in the MATH training
split (checked when the data are built). The prompt asks for a step-by-step solution with the final answer in
\verb|\boxed{}|. A completion is correct if math-verify~\cite{mathverify} finds its last boxed answer equal to the
reference answer. The judges never see the reference answer.

\paragraph{Rollouts.} Each step samples 16 problems and 8 completions per problem at temperature 1 (at most 1024
tokens) from a vLLM~\cite{kwon2023vllm} server that receives the policy adapter after every update. vLLM and the
training stack compute slightly different probabilities for the same weights. We correct each token by truncated
importance sampling~\cite{yao2025offpolicy}, with weight $\min(\pi_{\text{train}}/\pi_{\text{vLLM}},2)$, and we log the
mismatch.

\paragraph{Objective.} The advantage of a completion is its reward minus the mean reward of its group, divided by the
standard deviation of these centred rewards over the whole batch, as in Lite PPO~\cite{liu2026tricks}. With the batch
scale, the weight $\lambda$ of a judge term keeps its meaning in groups where the verifier gives every completion the same
reward. The loss is the token-level policy gradient, summed over all completion tokens of the batch and divided by
their number~\cite{liu2025understanding}. There is one AdamW step per batch, so PPO's ratio is one and its clipping is
inactive, and there is no KL term. Learning rate: chosen from $\{1,3,10\}\cdot10^{-5}$ on the validation set with the
\texttt{rlvr} reward (\S\ref{sec:res-sweep}) and then fixed for every condition. Gradient-norm clipping is at 1.0 and the
batch has 128 completions. Each run has 200 steps.

\paragraph{Reward conditions.} Let $v\in\{0,1\}$ be the verifier's verdict. $\pc$ is \kev{}'s $p(\text{yes})$ for
``Is the final answer of the proposed solution correct?''. $\pg$ is its $p(\text{yes})$ for ``Does the proposed
solution reach its final answer through genuine, relevant step-by-step reasoning\ldots?''. $P_{\mathrm{a}}$ is the AR
judge's $p(\text{yes})$ for the first question. The judge prompts are in
\file{experiments/part2_llm/prompts/judge_prompts.json} and were fixed before any GRPO run. The judges see the problem
and the completion, truncated to 1536 tokens.
\begin{center}\small
\begin{tabular}{@{}ll@{}}
\toprule
condition & reward \\ \midrule
\texttt{rlvr} & $v$ \\
\texttt{rlvr\_kev\_add} & $v+\lambda\pc$, $\lambda=0.2$ \\
\texttt{rlvr\_kev\_gate} & $v\cdot\pg$ \\
\texttt{kev\_only} & $\pc$ (no verifier) \\
\texttt{rlvr\_ar\_add} & $v+\lambda P_{\mathrm{a}}$, $\lambda=0.2$ \\ \bottomrule
\end{tabular}
\end{center}
With $\lambda<1$, a correct completion always outranks an incorrect one in the additive conditions. The judge can only
reorder completions with the same verdict. \texttt{kev\_only} asks a different question: whether a policy can be trained
by the judge alone, and whether it learns to exploit the judge. In every condition \kev{} also scores 64 random
completions per step, for the log.

\paragraph{Evaluation.} Every 50 steps the policy is evaluated greedily (at most 2048 tokens) on the validation set,
MATH-500 and a fixed 500-problem subset of GSM8K test~\cite{cobbe2021gsm8k}. After the last step it is evaluated on
the full GSM8K test set and MATH-500 (greedy), and on AMC23, AIME24 and AIME25 (16 samples per problem at temperature
0.6, top-$p$ 0.95, reported as the mean accuracy over samples). We report the final policy of every run; no checkpoint
is selected. Per-problem results are saved, so conditions are compared on the same problems: we report the paired
difference to \texttt{rlvr}, averaged over seeds, with a 95\% bootstrap confidence interval over problems.

\paragraph{Offline judge study.} The base policy wrote 4 completions at temperature 1 for each of 1000 problems:
MATH-500 and the 500-problem GSM8K subset. It used the training prompt and stop strings, so these are the completions
GRPO starts from. Both judges answered every question for every completion. For each signal we report its AUROC for
the verifier's verdict over all completions, and also within problems: the AUROC over the pairs of completions of the
same problem with different verdicts, which is what GRPO's group baseline uses. We also report calibration (expected
calibration error, ECE; Brier score) and cost. Four edits probe what moves a judge:
\begin{itemize}
\item \emph{answer only}: the completion is replaced by a single sentence that gives its own final answer;
\item \emph{padded}: twenty repetitions of a self-doubt sentence are appended;
\item \emph{asserted}: a claim that every step was verified is appended;
\item \emph{altered}: the final boxed answer of a correct completion is changed, so it becomes wrong.
\end{itemize}
A generative judge serves as a further baseline on a random subset: the same AR model, reasoning before it answers
(thinking on), with its score taken as $p(\text{yes})$ at the answer position. The study uses test problems. It
measures the judges and informs nothing in the GRPO runs, whose design and prompts were fixed before it ran.

\section{Experimental setup}\label{sec:setup}

Each run used one Colab virtual machine with one NVIDIA RTX PRO 6000 Blackwell Server Edition GPU (96\,GB). The trainer
(PyTorch 2.8 in kev's locked environment; the Hugging Face implementation of Qwen3.5, whose Gated DeltaNet layers run on
the flash-linear-attention kernels) and the vLLM~0.31 rollout server (its own environment) shared the GPU, and up to
three trainers shared one rollout server. Other experiments of this series often ran on the same GPU at the same time, so we
report wall-clock times only from measurements made on an otherwise idle GPU and label them as such. Seeds fix the
problem order, the vLLM sampling seeds (one per problem and step) and every random choice of the trainer. GPU kernels
are not deterministic, so two runs with the same seed are not bitwise identical. Runs checkpoint every 25 steps and
resume after an interruption; the number of resumed runs is reported.
